\section{Transaction ও ACID Properties}
\label{sec:ch06-transaction}
রহিম করিমকে 2000 টাকা পাঠাচ্ছে। রহিমের হিসাব থেকে টাকা কমল, কিন্তু করিমের হিসাবে যোগ হওয়ার আগেই কাজ থেমে গেল। শুধু প্রথম পরিবর্তন স্থায়ী হলে হিসাব ভুল হবে। তাই সম্পর্কিত কাজগুলোকে একটি সম্পূর্ণ কাজ হিসেবে পরিচালনা করতে হয়।
\begin{definitionbox}{Transaction (ট্রানজ্যাকশন)}
ডেটাবেজে একটি যৌক্তিক কাজ সম্পন্ন করতে এক বা একাধিক সংশ্লিষ্ট অপারেশন একসঙ্গে পরিচালনা করা হয়। এই কাজের একককে ট্রানজেকশন বলে।
\end{definitionbox}
একটি অর্ডারের পণ্য বিক্রি, মজুত কমানো ও পরিশোধের তথ্য রাখা সম্পর্কিত কাজ হতে পারে। কোন কাজগুলো একই ট্রানজেকশনে থাকবে তা ব্যবস্থার নিয়ম নির্ধারণ করে। কেবল একাধিক SQL লিখলেই সেগুলো স্বয়ংক্রিয়ভাবে একটি সম্পূর্ণ লেনদেন হয় না; শুরু ও সফল সমাপ্তি বা বাতিলের সীমানাও নির্ধারণ করতে হয়।

\subsubsection{Transaction Boundary}
একটি Transaction-এ একটিমাত্র SQL নির্দেশও থাকতে পারে; ভর্তি নিশ্চিত করা বা অর্ডারের মজুত হালনাগাদের মতো কাজের সীমানা প্রয়োজন অনুযায়ী নির্ধারিত হয়।

অনেক ব্যবস্থায় একটি স্বতন্ত্র নির্দেশ নিজস্ব ট্রানজেকশনে সম্পন্ন হয়। কয়েকটি নির্দেশ একসঙ্গে সফল বা বাতিল করতে হলে সেগুলোকে একই ট্রানজেকশনে রাখতে হয়। টাকা কমানোর নির্দেশটি আগে COMMIT করা থাকলে পরে শুরু করা অন্য ট্রানজেকশনের ROLLBACK সেই পরিবর্তন ফিরিয়ে দেবে না। নিচের উদাহরণে তাই কাজের শুরু ও সমাপ্তির সীমানা গুরুত্বপূর্ণ।


\subsection{একই লেনদেনের আগে ও পরে}
এখানে দুই হিসাবই আগে থেকে আছে, রহিমের পর্যাপ্ত টাকা আছে এবং কোনো ফি কাটা হচ্ছে না। AccountID প্রতিটি হিসাবের অনন্য পরিচিতি; টাকার পরিমাণ পূর্ণ টাকা হিসেবে দেখানো হয়েছে।
\begin{chSixSideBySide}[.48]
\small\textbf{Account: লেনদেনের আগে}\par
\begin{tabular}{|l|l|l|}\hline\chSixHeader
AccountID & Name & Balance\\\hline
A01 & Rahim & 10000\\\hline
A02 & Karim & 3000\\\hline
\end{tabular}
\chSixNextColumn
\small\textbf{সফল COMMIT-এর পরে}\par
\begin{tabular}{|l|l|l|}\hline\chSixHeader
AccountID & Name & Balance\\\hline
A01 & Rahim & 8000\\\hline
A02 & Karim & 5000\\\hline
\end{tabular}
\end{chSixSideBySide}
শুধু প্রথম ধাপ হলে সাময়িকভাবে 8000 ও 3000 দেখা যেতে পারে; সেটিকে সম্পূর্ণ লেনদেনের চূড়ান্ত ফল হিসেবে গ্রহণ করা যাবে না। সফল হলে দুই পরিবর্তনই থাকবে; বাতিল হলে এই লেনদেনের দুই পরিবর্তনই ফিরবে।


\textbf{ব্যাংকের এই সমস্যায় চারটি আলাদা নিশ্চয়তা দরকার।} এক ধাপের ফল স্থায়ী রেখে থেমে যাওয়া চলবে না; হিসাবের বৈধতার নিয়ম ভাঙা যাবে না; একই সময়ে অন্য কাজ হলে অসংগত ফল ঠেকাতে হবে; সফলভাবে গৃহীত ফল পুনরায় চালুর পরেও বহাল থাকতে হবে। এগুলো একই কথার চারটি নাম নয়। ACID-এর প্রতিটি বৈশিষ্ট্য নিচের উদাহরণে ভিন্ন ঝুঁকি সামলায়।

\subsection{ACID-এর চারটি বৈশিষ্ট্য}
ACID হলো Atomicity, Consistency, Isolation ও Durability-এর আদ্যক্ষর। চারটি বৈশিষ্ট্য একই সমস্যার চারটি আলাদা দিক বোঝায়।
\begin{enumerate}
\item \textbf{Atomicity:} লেনদেনের সব পরিবর্তন সম্পন্ন হবে অথবা কোনোটিই স্থায়ী হবে না। Debit সফল কিন্তু Credit ব্যর্থ হলে অসম্পূর্ণ লেনদেন বাতিল করতে হয়।
\item \textbf{Consistency:} লেনদেনের আগে ও পরে নির্ধারিত শুদ্ধতার নিয়ম বজায় থাকবে। বৈধ হিসাব, অনুমোদিত Balance এবং ফিবিহীন স্থানান্তরে মোট টাকা অক্ষুণ্ণ থাকতে হয়।
\item \textbf{Isolation:} সমসাময়িক লেনদেনের পারস্পরিক প্রভাব নিয়ন্ত্রণ করে সঠিক ফল বজায় রাখা হয়। একই Balance নিয়ে দুটি উত্তোলনে অসামঞ্জস্য প্রতিরোধ করতে হয়।
\item \textbf{Durability:} সফল COMMIT-এর পরিবর্তন স্থায়ী সংরক্ষণে বজায় থাকবে। ব্যবস্থা হঠাৎ বন্ধ হওয়ার পর স্বাভাবিক পুনরুদ্ধারে COMMIT করা স্থানান্তর বহাল থাকে।
\end{enumerate}
\textbf{Consistency বুঝতে হবে নিয়ম দিয়ে।} এই উদাহরণে মোট টাকা $10000+3000=8000+5000=13000$। কিন্তু বাইরের উৎস থেকে নতুন টাকা জমা হলে মোট টাকা বাড়াই স্বাভাবিক। তাই সব ট্রানজেকশনে মোট টাকা অপরিবর্তিত থাকা Consistency-এর সংজ্ঞা নয়। ঋণাত্মক Balance নিষিদ্ধ হলে সেই বিধি মানাও এর একটি দিক। সঠিক অ্যাপ্লিকেশন-যুক্তি ও ডেটাবেজে নির্ধারিত বিধি মিলিয়ে নিয়ম কার্যকর করতে হয়।

\textbf{Isolation বুঝতে হবে একই সময়ের কাজ দিয়ে।} ধরা যাক 10000 টাকা থাকা হিসাব থেকে দুজন একই সময়ে 7000 করে তুলতে চায়। উভয়ে পুরোনো Balance দেখে অনুমোদন পেলে সমস্যা হবে। ব্যবস্থাকে সংঘাত শনাক্ত করে অপেক্ষা, পুনরায় পরীক্ষা বা একটি কাজ বাতিলের মাধ্যমে নিয়ম রক্ষা করতে হয়। Isolation মানে সব ব্যবহারকারীকে সবসময় একে একে কাজ করতেই হবে—এমন নয়।

\textbf{Durability ও ব্যাকআপ আলাদা।} COMMIT-এর পরে বিদ্যুৎ চলে গেলেও গৃহীত লেনদেনের ফল ধরে রাখার ব্যবস্থা DBMS-এ থাকে। তবে সব সংরক্ষণযন্ত্র ও নিরাপদ কপি একসঙ্গে ধ্বংস হলেও তথ্য কখনো হারাবে না—Durability এমন নিশ্চয়তা দেয় না। তাই ব্যাকআপ ও ডেটা পুনরুদ্ধারের ব্যবস্থাও প্রয়োজন।

\subsection{BEGIN, COMMIT ও ROLLBACK-এর প্রয়োগ}
নিচের দুটি উদাহরণ SQLite-এর জন্য। প্রতিটি উদাহরণ স্বাধীনভাবে আগের 10000 ও 3000 Balance থেকে শুরু হয়েছে। প্রথমটিতে সফল স্থানান্তর স্থায়ী করা হয়েছে।
\begin{chSixSideBySide}[.60]
% transaction-example: commit
\begin{verbatim}
BEGIN TRANSACTION;
UPDATE Account
SET Balance = Balance - 2000
WHERE AccountID = 'A01';
UPDATE Account
SET Balance = Balance + 2000
WHERE AccountID = 'A02';
COMMIT;
SELECT AccountID, Balance
FROM Account ORDER BY AccountID;
\end{verbatim}
\chSixNextColumn
\small\textbf{ফল}\par
\begin{tabular}{|l|l|}\hline\chSixHeader
AccountID & Balance\\\hline
A01 & 8000\\\hline A02 & 5000\\\hline
\end{tabular}\par\smallskip
BEGIN লেনদেন শুরু করে; COMMIT দুই পরিবর্তন স্থায়ী করে। WHERE দিয়ে প্রতিবার নির্দিষ্ট হিসাবটি বেছে নেওয়া হয়েছে।
\end{chSixSideBySide}
\begin{chSixSideBySide}[.60]
% transaction-example: rollback
\begin{verbatim}
BEGIN TRANSACTION;
UPDATE Account
SET Balance = Balance - 2000
WHERE AccountID = 'A01';
ROLLBACK;
SELECT AccountID, Balance
FROM Account ORDER BY AccountID;
\end{verbatim}
\chSixNextColumn
\small\textbf{বাতিলের পর ফল}\par
\begin{tabular}{|l|l|}\hline\chSixHeader
AccountID & Balance\\\hline
A01 & 10000\\\hline A02 & 3000\\\hline
\end{tabular}\par\smallskip
দ্বিতীয় ধাপে না গিয়ে বাতিলের সিদ্ধান্ত নেওয়া হয়েছে। ROLLBACK প্রথম পরিবর্তনও ফিরিয়েছে; টাকা শুধু এক হিসাব থেকে হারিয়ে যায়নি।
\end{chSixSideBySide}
এটি লেনদেন বোঝার সরল উদাহরণ। বাস্তব ব্যবস্থায় দুই হিসাবের অস্তিত্ব, পর্যাপ্ত অর্থ, অনুমতি এবং প্রতিটি UPDATE-এর ফল পরীক্ষা করে তবেই COMMIT করতে হয়। ভুল AccountID-তে কোনো সারি না বদলালেও UPDATE সবসময় ত্রুটি দেখায় না। আবার কোনো নির্দেশ ব্যর্থ হলেই SQLite সব পূর্ববর্তী নির্দেশ নিজে থেকে বাতিল করবে—এমনও ধরে নেওয়া যাবে না; ব্যর্থতা সামলানোর যুক্তি প্রয়োজন।


\begin{flushleft}\small
\begin{tabularx}{\linewidth}{|p{24mm}|Y|Y|}\hline\chSixHeader
বিবেচনা & COMMIT & ROLLBACK\\\hline
উদ্দেশ্য & চলমান লেনদেনের পরিবর্তন স্থায়ী করা & চলমান লেনদেনের পরিবর্তন বাতিল করা\\\hline
এই উদাহরণে & দুই ধাপের পরে 8000 ও 5000 বহাল & বাতিল হলে আবার 10000 ও 3000\\\hline
সিদ্ধান্ত & প্রয়োজনীয় কাজ ও যাচাই সফল & ত্রুটি বা বাতিলের সিদ্ধান্ত\\\hline
\end{tabularx}
\end{flushleft}
\subsubsection{Transaction-এর ব্যবহার}
\begin{enumerate}
\item \textbf{টিকিট বরাদ্দ:} টিকিট বরাদ্দের সঙ্গে আসনের অবস্থা বদলানোর কাজ এক লেনদেনে রাখা যায়।
\item \textbf{বিক্রয় ও মজুত:} বিক্রয়ের সঙ্গে মজুত হালনাগাদের কয়েকটি কাজ এক লেনদেনে রাখা যায়।
\end{enumerate}
\subsubsection{সীমাবদ্ধতা}
\begin{enumerate}
\item \textbf{ব্যবস্থাপনার খরচ:} সমসাময়িক কাজ নিয়ন্ত্রণে অপেক্ষা ও অতিরিক্ত ব্যবস্থাপনার খরচ হতে পারে।
\item \textbf{সঠিক নিয়মের প্রয়োজন:} ভুল নিয়ম লিখলে ট্রানজেকশন সেই নিয়মকে নিজে সঠিক করে না; কোন কাজগুলো একসঙ্গে সফল হতে হবে তা আগে নির্ধারণ করতে হয়।
\end{enumerate}

\subsection{Transaction, COMMIT ও ROLLBACK-এর পার্থক্য}
\begin{flushleft}\small
\begin{tabularx}{\linewidth}{|p{23mm}|Y|Y|Y|}\hline\chSixHeader
বিবেচনা & Transaction & COMMIT & ROLLBACK\\\hline
ভূমিকা & একসঙ্গে সম্পন্ন হওয়ার জন্য নির্ধারিত কাজের একক & সফল লেনদেনের পরিবর্তন স্থায়ী করে & চলমান লেনদেনের পরিবর্তন বাতিল করে\\\hline
সীমা & কোন কাজগুলো একসঙ্গে সফল হবে তা নির্ধারিত থাকে & আগে COMMIT করা ভুলের জন্য নতুন অনুমোদিত পরিবর্তন বা উপযুক্ত Recovery দরকার & পুরোনো COMMIT করা পরিবর্তন বাতিল করার নির্দেশ নয়\\\hline
\end{tabularx}
\end{flushleft}
\textbf{SAVEPOINT:} লেনদেনের মাঝখানে ফেরার চিহ্ন; পুরো লেনদেন বাতিল না করেও নির্দিষ্ট অংশ ফেরানোর কাজে লাগে।
\begin{enumerate}
\item \textbf{Atomicity:} কাজের সম্পূর্ণতা বোঝায়।
\item \textbf{Consistency:} নিয়মের শুদ্ধতা বোঝায়।
\item \textbf{Isolation:} সমসাময়িক কাজের প্রভাব নিয়ন্ত্রণ বোঝায়।
\item \textbf{Durability:} COMMIT করা ফলের স্থায়িত্ব বোঝায়।
\end{enumerate}

\subsubsection{আংশিক স্থানান্তরে মোট টাকার পরিবর্তন}
উদাহরণে দুই হিসাবের মোট টাকা 13000। শুধু 2000 টাকা কমানোর ধাপ স্থায়ী হলে মোট 11000 হয়ে যাবে, যা এই ফিবিহীন স্থানান্তরের নিয়ম ভাঙে। সফল COMMIT-এ স্থিতি 8000 ও 5000; বাতিল হলে 10000 ও 3000 ফিরে আসে।
